\begin{abstract}
We introduce \textsc{QLoRA}, a memory-efficient finetuning method that enables the finetuning of large language models with up to 65 billion parameters on a single 48GB GPU, while maintaining the performance typical of full 16-bit finetuning. \textsc{QLoRA} operates by propagating gradients through a static, 4-bit quantized pretrained language model into Low Rank Adapters~(LoRA). Our top-performing model family, named \textbf{Guanaco}, surpasses all previously available open-source models on the Vicuna benchmark, achieving 99.3\% of ChatGPT's performance after just 24 hours of finetuning on a single GPU. \textsc{QLoRA} introduces several novel strategies to minimize memory usage while retaining high performance: (a) 4-bit NormalFloat (NF4), a new numerical format designed to be information-theoretically optimal for normally distributed weights, (b) Double Quantization, which reduces memory usage by further quantizing the quantization coefficients, and (c) Paged Optimizers, designed to handle memory spikes effectively. Leveraging \textsc{QLoRA}, we finetuned over 1,000 models, offering an extensive evaluation of instruction following and chatbot capabilities across eight instruction datasets, different model architectures (LLaMA, T5), and various model sizes, including those (e.g., 33B and 65B) not practical to finetune conventionally. Our findings demonstrate that applying QLoRA finetuning to small, high-quality datasets delivers state-of-the-art results, even with smaller models than previous leading methods. We conduct an in-depth assessment of chatbot performance utilizing both human annotators and GPT-4 as evaluators, showing GPT-4 serves as a cost-effective and reliable alternative. Additionally, we highlight that current benchmarks for chatbots may not provide reliable measurements of real performance. Through targeted failure case analysis, we illustrate the specific shortcomings of \textbf{Guanaco} when compared with ChatGPT. All models and associated code are released, including CUDA kernels optimized for 4-bit training.\footnote{\url{https://github.com/artidoro/qlora} and \url{https://github.com/TimDettmers/bitsandbytes}}
\end{abstract}

\section{Introduction}

Refining large language models (LLMs) through finetuning has proven to be a powerful strategy for enhancing model capabilities \citep{min2021metaicl, wei2021finetuned, ouyang2022training, wang2022super, wang2022self, liu2022few}, as well as for incorporating or removing specific behaviors \citep{ouyang2022training,askell2021general,bai2022training}. Nevertheless, the process of finetuning models at this scale incurs substantial computational costs; for instance, conventional 16-bit finetuning of a LLaMA model with 65B parameters~\citep{touvron2023llama} demands over 780 GB of GPU memory. Although recent advances in quantization can diminish the inference memory requirements for LLMs \citep{dettmers2022llm, dettmers2022case, frantar2022gptq, xiao2022smoothquant}, these approaches typically fail to support efficient training, as evidenced by their instability during the learning process \citep{wortsman2023stable}.

In this work, we introduce a novel method for finetuning LLMs quantized to 4 bits, achieving this without any sacrifice in performance. Our approach, \textsc{QLoRA}, leverages a new high-precision quantization strategy to compress pretrained models to 4-bit precision and integrates a compact set of trainable Low-rank Adapter weights \citep{hu2021lora}. 

Finetuning is performed by propagating gradients through the quantized model parameters, with updates restricted to these adapter weights.

Employing \textsc{QLoRA}, we reduce the GPU memory required to finetune a 65B parameter model from in excess of 780GB to less than 48GB, while maintaining runtime efficiency and predictive accuracy relative to traditional 16-bit full finetuning. This development dramatically expands the feasibility of LLM finetuning: models of unprecedented size now become accessible for finetuning on a single GPU. Utilizing our method, we developed the \textbf{Guanaco} model suite. The second strongest \textbf{Guanaco} variant attains 97.8\% of ChatGPT’s performance on the Vicuna benchmark~\citep{vicuna2023}, while training can be completed in under 12 hours on a consumer GPU; employing a single professional GPU for 24 hours, our largest variant closes the performance gap further, reaching 99.3\% of ChatGPT’s Vicuna score. Notably, in deployment, our smallest \textbf{Guanaco} model, which has 7B parameters, fits into just 5 GB of memory and surpasses the 26 GB Alpaca model by over 20 percentage points on the Vicuna benchmark (see Table~\ref{tbl:vicuna_eval}).

\textsc{QLoRA} implements several key advancements to decrease memory requirements while maintaining model quality: (1) \textbf{4-bit NormalFloat}, an information-theoretically optimal quantization format tailored for data following a normal distribution, which demonstrates superior empirical outcomes relative to conventional 4-bit Integers and 4-bit Floats. (2) \textbf{Double Quantization}, a technique in which the quantization coefficients themselves are quantized, leading to an average reduction of around 0.37 bits per parameter (translating to about 3 GB saved for a 65B-parameter model). (3) \textbf{Paged Optimizers}, which leverage NVIDIA unified memory to circumvent the memory spikes caused by gradient checkpointing during the processing of mini-batches with extensive sequence lengths. By integrating these advances, we present an enhanced LoRA methodology that applies adapters throughout every layer of the network, thereby circumventing most of the accuracy compromises encountered in preceding research.

The improved memory efficiency provided by \textsc{QLoRA} allows for an extensive exploration of instruction fine-tuning and chatbot capabilities at model scales that would be unattainable with standard fine-tuning methods due to prohibitive memory consumption. Consequently, we trained over 1,000 models on a diverse array of instruction tuning datasets, model architectures, and sizes ranging from 80M to 65B parameters. Beyond demonstrating that \textsc{QLoRA} matches 16-bit performance (\S\ref{sec:comp_to_std}) and enables the training of a state-of-the-art chatbot, \textbf{Guanaco} (\S\ref{sec:pushing}), we also examine model trends across the training runs. Notably, our results indicate that the quality of training data surpasses dataset size in importance—for instance, a compact 9k-example dataset (OASST1) outperformed a significantly larger 450k-example dataset (FLAN v2, subsampled) with respect to chatbot effectiveness, despite both being aimed at enhancing instruction-following abilities. Furthermore, we find that exceptional results on the Massive Multitask Language Understanding (MMLU) benchmark do not necessarily correlate with high Vicuna chatbot benchmark performance and vice versa, emphasizing that alignment between dataset characteristics and target tasks is more critical than raw dataset magnitude.

In addition, we conduct a comprehensive study evaluating chatbot performance through assessments by both human annotators and GPT-4. For this evaluation, we implement a tournament-style framework wherein models engage in direct matchups, striving to generate the most effective response for given prompts. The outcome of each match is determined either by human judges or by GPT-4, and aggregate results from these matchups are used to compute Elo scores~\citep{elo1967proposed, elo1978rating}, which serve as a basis for ranking the chatbots. Our findings indicate substantial alignment between the rankings assigned by GPT-4 and those by humans, though there are notable cases where their judgments diverge. Consequently, we emphasize that while automated, model-based evaluations offer a more resource-efficient alternative to human annotation, they also come with inherent uncertainties.

To supplement these benchmarking results, we perform an in-depth qualitative analysis of the \textbf{Guanaco} models. This exploration brings to light various strengths and shortcomings of the models that are not fully revealed through quantitative metrics alone.

For reproducibility and further research, we publicly release all model outputs with accompanying human and GPT-4 ratings. Our codebase, along with custom CUDA kernels, is open-sourced and our methods have been integrated into the Hugging Face transformers library~\citep{wolf2019huggingface}, ensuring broad accessibility. Furthermore, we provide adapter checkpoints for 7/13/33/65B model variants, each finetuned on eight separate instruction-following datasets, yielding a total of 32 distinct, openly available finetuned models.

\section{Background}
\paragraph{Block-wise k-bit Quantization} Quantization involves transforming data from a high-precision format to a lower-precision one, typically by mapping data types with greater bit-width (such as 32-bit floating points) to those with fewer bits (like 8-bit integers). This process often requires adjusting the original data so that the reduced-bit representation covers the full target range, commonly by scaling the input tensor with respect to its absolute maximum value. For instance, the quantization of a 32-bit floating point (FP32) tensor into an Int8 tensor, whose valid values range from $[-127, 127]$, is expressed as:

\begin{equation}
    \mathbf{X}^{ \text{Int8}} = \text{round}\left( \frac{127}{{\text{absmax}}(\mathbf{X}^{{\text{FP32}}})}\mathbf{X}^{{\text{FP32}}} \right) = \text{round}(c^{\text{FP32}}\cdot \mathbf{X}^{{\text{FP32}}}),
\end{equation}

where $c$ denotes the {\em quantization constant} or {\em scale factor}. The corresponding dequantization operation, which reconstructs the higher-precision values from their quantized counterparts, is defined as:

\begin{equation}
    \text{dequant}(c^{\text{FP32}}, \mathbf{X}^{\text{Int8}}) = \frac{\mathbf{X}^{\text{Int8}}}{c^{\text{FP32}}} = \mathbf{X}^{\text{FP32}}
\end{equation}

A limitation of this method is that when the input tensor contains values with large magnitudes—i.e., outliers—quantization bins, corresponding to specific bit patterns, become inefficiently populated, with some bins rarely or never being used. To address this outlier problem, a widely adopted strategy involves partitioning the input tensor into smaller segments, or blocks, which are then quantized separately. Each block is assigned its own quantization constant $c$. More precisely, the input tensor $\mathbf{X}\in \mathbb{R}^{b\times h}$ is first flattened and then divided into $n$ consecutive blocks, each of length $B$, yielding ${n = ({b\times h} )/{B} }$ blocks in total. Each of these blocks is quantized independently, as described in Equation~1, resulting in a quantized representation as well as $n$ distinct quantization constants $c_i$.

\paragraph{Low-rank Adapters} LoRA finetuning \citep{hu2021lora} introduces a parameter-efficient approach by incorporating a limited number of trainable parameters, referred to as adapters, and keeping the remainder of the model parameters unchanged. During stochastic gradient descent, gradients are routed through the immutable pretrained weights and are used exclusively to update the adapter modules to minimize the loss. LoRA modifies a standard linear transformation by appending an extra low-rank, factorized term. Specifically, for a projection expressed as ${\mathbf{X} \mathbf{W} = \mathbf{Y}}$ with $\mathbf{X}\in  \mathbb{R}^{b\times h}$ and $\mathbf{W}\in  \mathbb{R}^{h\times o}$, the LoRA formulation is:

\begin{equation}
    \mathbf{Y} = \mathbf{X}\mathbf{W} + s\mathbf{X}\mathbf{L}_1\mathbf{L}_2,
\end{equation}

with $\mathbf{L}_1\in  \mathbb{R}^{h\times r}$, $\mathbf{L}_2\in  \mathbb{R}^{r\times o}$, and $s$ denoting a scalar scaling factor.

\paragraph{Memory Requirement of Parameter-Efficient Finetuning} A key aspect to consider is the memory consumption associated with LoRA during training, specifically regarding both the quantity and size of adapters implemented. Owing to LoRA's very small memory usage, it is feasible to utilize additional adapters to enhance model effectiveness with only marginal impact on overall memory demands. Although LoRA is introduced as a Parameter Efficient Finetuning (PEFT) approach, it is important to note that, for LLM finetuning, the principal memory cost arises from storing activation gradients rather than the LoRA parameters themselves. For example, in the case of training a 7B LLaMA model on FLAN v2 with a batch size of 1 and LoRA parameters set to the standard 0.2\% of the base model size \citep{hu2021lora,liu2022few}, input gradients for LoRA account for 567 MB of memory usage, compared to just 26 MB for the LoRA weights. Utilizing gradient checkpointing~\citep{chen2016training} can reduce these input gradients to roughly 18 MB per sequence—still more than the total memory taken up by all LoRA weights. For reference, the memory required by the 4-bit base model is 5,048 MB. This illustrates not only the significance of gradient checkpointing in managing memory costs, but also demonstrates that further decreasing LoRA parameter size leads to limited memory savings. Consequently, there is considerable flexibility to increase the number of adapters with little effect on the aggregate memory usage during training (refer to Appendix~\ref{app:memory} for detailed information). As explained subsequently, this flexibility is important for achieving performance equivalent to full 16-bit precision models.

\section{\textsc{QLoRA} Finetuning}

\textsc{QLoRA} enables accurate 4-bit finetuning by employing two principal techniques that we introduce: 4-bit NormalFloat (NF4) quantization and Double Quantization. To further address memory inefficiencies during gradient checkpointing—commonly leading to out-of-memory errors and impeding single-machine finetuning of large-scale models—we propose Paged Optimizers.

In \textsc{QLoRA}, one utilizes a low-precision storage format, typically 4-bit, along with a separate computation data type, most often BFloat16. In practical terms, when a \textsc{QLoRA} weight tensor is needed for computation, it is first dequantized to BFloat16, after which matrix multiplication proceeds in 16-bit precision.

Below, we elaborate on each of the foundational techniques of \textsc{QLoRA} before formally defining the approach.

\paragraph{4-bit NormalFloat Quantization} The NormalFloat (NF) format extends Quantile Quantization as described in \citep{dettmers2022optimizers}. This quantization approach is theoretically optimal in terms of information preservation, as it allocates the same number of input tensor values to each quantization bin, thereby equalizing representation. Quantile quantization functions by using the empirical cumulative distribution function to estimate quantiles from the input tensor.

A primary drawback of quantile quantization lies in its computational cost, as accurately estimating quantiles is resource-intensive. For this reason, efficient quantile approximation strategies like SRAM quantiles~\citep{dettmers2022optimizers} are employed. Nevertheless, the inherent approximation in these algorithms can cause large quantization errors, especially for outlier values—which are often critical.

These issues, namely costly quantile computations and approximation inaccuracies, can be circumvented when input tensors are sampled from distributions invariant except for a quantization scaling factor. In such settings, input tensors share identical quantiles, which renders exact quantile calculation much more tractable.

Because pretrained neural network weights are typically distributed according to a zero-mean normal law with standard deviation $\sigma$ (refer to Appendix~\ref{app:normality}), we can map all such weights onto a unified reference distribution by adjusting $\sigma$ so that the resulting distribution precisely spans the allowable range of our chosen data type. For this study, we designate the data type interval to be $[-1, 1]$. Accordingly, both the data type quantiles and the neural network weight values must be scaled to fall within this interval.

The information-theoretically ideal data type for representing zero-mean Gaussian distributions with arbitrary standard deviation $\sigma$ over the domain $[-1, 1]$ can be constructed using the following steps: (1) Compute the $2^k+1$ quantile boundaries from the standard normal distribution $N(0, 1)$ to form a $k$-bit quantile-quantized data type suitable for Gaussian-distributed values, (2) Normalize these data type values so that they lie within the $[-1, 1]$ interval, and (3) Quantize the input weight tensor by likewise normalizing it into the $[-1, 1]$ range, typically by dividing by its largest absolute value.

Once the domains of the weight tensor and the data type coincide, conventional quantization can be applied. The third step effectively rescales the weight tensor’s standard deviation to equal that of the $k$-bit data type. More precisely, the $2^k$ data type values $q_i$ can be defined as:

\begin{equation}
q_i  = \frac{1}{2}\left( Q_X\left(\frac{i}{2^k+1}\right) + Q_X\left(\frac{i+1}{2^k+1}\right)\right),
\end{equation}

where $Q_X(\cdot)$ denotes the quantile function for the standard normal distribution $N(0,1)$. A limitation of symmetric $k$-bit quantization is its inability to precisely represent the value zero—an essential characteristic needed for quantizing padding and other elements with zero values without introducing error. To address this and guarantee that zero maps exactly to a quantized value, while utilizing all $2^k$ representable codes for a $k$-bit data type, we propose an asymmetric variant. Specifically, we estimate quantiles $q_i$ over two intervals: $2^{k-1}$ quantiles for the negative domain and $2^{k-1}+1$ for the positive, then merge these two sets of $q_i$ and eliminate one of the two zeros shared by both sets. We refer to the resulting data type, in which each quantization bin is expected to contain an equal proportion of the data, as \emph{k-bit NormalFloat} (NFk). This scheme is information-theoretically optimal for data that follow a zero-centered normal distribution. Details of the precise values for this data type are presented in Appendix~\ref{app:nf4}.

\paragraph{Double Quantization{}}
We propose a method called \emph{Double Quantization} (DQ), wherein quantization constants themselves are further quantized to achieve additional reductions in memory usage. While employing small blocksizes is necessary for high-precision 4-bit quantization~\citep{dettmers2022case}, it leads to non-negligible memory costs. For instance, with a blocksize of 64 and 32-bit quantization constants for $\mathbf{W}$, these constants require on average $32/64 = 0.5$ bits per parameter. Double Quantization alleviates this overhead by compressing the storage requirements for the quantization constants.

Concretely, in Double Quantization, the quantization constants $c_2^{\text{FP32}}$ used in the initial quantization stage are treated as inputs to a subsequent quantization process. This second stage produces quantized constants $c_2^{\text{FP8}}$ and a new set of quantization constants $c_1^{\text{FP32}}$. We utilize 8-bit floating-point representations and a blocksize of 256 for this second quantization phase, as previous findings~\citep{dettmers2022case} confirm that 8-bit quantization maintains performance. Since $c_2^{\text{FP32}}$ are always positive, we subtract their mean to recenter the distribution at zero prior to quantization, enabling the use of symmetric quantization. On average, for a blocksize of 64, this technique reduces the memory demand per parameter from $0.5$ bits ($32/64$) to $8/64 + 32/(64\cdot256) = 0.127$ bits, achieving a reduction of 0.373 bits per parameter.

\paragraph{Paged Optimizers} leverage the NVIDIA unified memory~\footnote{\tiny \url{https://docs.nvidia.com/cuda/cuda-c-programming-guide}} capability, which transparently manages page-level data transfers between the CPU and GPU. This mechanism ensures seamless GPU computations, particularly in situations where the GPU exhausts its memory resources. Unified memory functions analogously to traditional CPU RAM-to-disk paging systems. We utilize this approach to reserve paged memory for optimizer states; when GPU memory is insufficient, these states are offloaded automatically to CPU RAM and subsequently reloaded into GPU memory during the optimizer's update phase as required.

\paragraph{\textsc{QLoRA}.} Building upon the aforementioned components, we formulate \textsc{QLoRA} for a quantized single linear layer incorporating a single LoRA adapter as follows:

\begin{equation}
    \mathbf{Y}^{\text{BF16}} =  \mathbf{X}^{\text{BF16}} \text{doubleDequant}(c_1^{\text{FP32}}, c_2^{\text{k-bit}}, \mathbf{W}^{\text{NF4}}) + \mathbf{X}^{\text{BF16}}\mathbf{L}^{\text{BF16}}_1\mathbf{L}^{\text{BF16}}_2,
\end{equation}

with doubleDequant$(\cdot)$ specified as:

\begin{equation}
    \text{doubleDequant}(c_1^{\text{FP32}}, c_2^{\text{k-bit}}, \mathbf{W}^{\text{k-bit}}) = \text{dequant}(\text{dequant}(c_1^{\text{FP32}}, c_2^{\text{k-bit}}), \mathbf{W}^{\text{4bit}}) = \mathbf{W}^{\text{BF16}},
\end{equation}

In our implementation, $\mathbf{W}$ employs the NF4 format, while $c_2$ utilizes FP8.

To maximize quantization accuracy, a blocksize of 64 is selected for $\mathbf{W}$, whereas $c_2$ adopts a blocksize of 256 in order to optimize memory usage.

During parameter optimization, gradients are computed only with respect to the adapter weights, denoted as $\frac{\partial E}{\partial \mathbf{L}_i}$; gradients for the 4-bit weight matrix $\frac{\partial E}{\partial \mathbf{W}}$ are not required. Nonetheless, determining $\frac{\partial E}{\partial \mathbf{L}_i}$ necessitates calculating $\frac{\partial \mathbf{X}}{\partial \mathbf{W}}$, which follows from equation~(5) by dequantizing $\mathbf{W}^{\text{NF4}}$ to the computation format $\mathbf{W}^{\text{BF16}}$, thereby allowing the derivative $\frac{\partial \mathbf{X}}{\partial \mathbf{W}}$ to be evaluated with BFloat16 precision.

In summary, \textsc{QLoRA} employs a single data type for storage—commonly 4-bit NormalFloat—and a separate data type, 16-bit BrainFloat, for computation. During training, the stored data is dequantized into the computational format to enable both forward and backward passes. However, the computation of weight gradients is restricted to the LoRA parameters, which are represented in 16-bit BrainFloat.

\section{QLoRA vs. Standard Finetuning}
\label{sec:comp_to_std}

Our earlier discussion outlined the workings of QLoRA and its potential for substantial reductions in memory consumption during model finetuning. This naturally raises the question of whether QLoRA’s performance can match that of conventional full-parameter finetuning. We are also interested in dissecting QLoRA’s elements, with particular attention to the role that NormalFloat4 plays relative to traditional Float4. The experiments described in the following sections were designed to address these topics.

\paragraph{Experimental setup.} Our analysis spans three model architectures—encoder, encoder-decoder, and decoder-only—to facilitate comparisons between QLoRA, 16-bit adapter-based finetuning, and standard full-parameter finetuning on models scaling up to 3B parameters. The evaluations utilize GLUE \citep{wang2018glue} with RoBERTa-large \citep{liu2019roberta}, Super-NaturalInstructions (TKInstruct) \citep{wang2022super} with T5 \citep{t5}, and 5-shot MMLU \citep{hendrycksmeasuring}, where LLaMA is finetuned on Flan v2 \citep{longpre2023flan} and Alpaca \citep{alpaca}. To further investigate the merits of NF4 compared to alternative 4-bit formats, we adopt the evaluation strategy from \citet{dettmers2022case}, recording post-quantization zero-shot accuracy and perplexity across diverse models (OPT~\citep{zhang2022opt}, LLaMA~\citep{touvron2023llama}, BLOOM~\citep{scao2022bloom}, and Pythia~\citep{biderman2023pythia}), and varying model sizes from 125m up to 13B.

Further specifics regarding each experimental setup are provided alongside the corresponding results for clarity. The appendix (Appendix~\ref{app:exp-setup}) contains comprehensive methodological details.

Paged optimizers enable \textsc{QLoRA} finetuning on models with 33B/65B parameters using a single GPU with 24/48GB memory. However, we refrain from supplying quantitative measurements for Paged Optimizers, given that paging is triggered primarily during training with unusually long sequences—a rare event. Nonetheless, we assess the training time implications of paged optimizers on 65B-parameter models with 48GB GPUs, observing that for a batch size of 16, training speed is indistinguishable from that of conventional optimizers. It would be worthwhile for future research to systematically identify and characterize scenarios where paging leads to slower training.

\paragraph{Default LoRA hyperparameters do not match 16-bit performance}

Applying the default LoRA configuration—targeting the query and value matrices of the attention mechanism as in \citep{hu2021lora}—does not enable reproduction of full-parameter finetuning outcomes for large-scale models. Figure~\ref{fig:hyper} illustrates that, in the case of LLaMA 7B finetuned on Alpaca, the dominant LoRA hyperparameter influencing performance is the total number of LoRA adapters used. Achieving parity with full-model finetuning necessitates the integration of LoRA adapters into all linear layers within the transformer block. Other LoRA configuration choices, such as the projection rank $r$, exhibit negligible impact on final performance (refer to Appendix \ref{app:exp-setup}).

In a similar vein, our investigation reveals that the standard hyperparameters used for complete finetuning baselines are insufficiently optimized. To establish strong baselines, we perform a systematic search across learning rates ranging from 1e-6 to 5e-5 and batch sizes from 8 to 128. The finetuning outcomes for the 7B LLaMA model on the Alpaca dataset are depicted in Figure~\ref{fig:hyper}.

\paragraph{4-bit NormalFloat Surpasses 4-bit Floating Point in Performance} 

Although the 4-bit NormalFloat (NF4) data type is known to be theoretically optimal from an information perspective, it remains uncertain whether these theoretical strengths confer practical benefits. To examine this, we replicate the methodology outlined in \citet{dettmers2022case}, evaluating quantized LLMs—specifically OPT \citep{zhang2022opt}, BLOOM \cite{scao2022bloom}, Pythia \cite{biderman2023pythia}, and LLaMA—at scales from 125M to 65B parameters using various data representations. Assessments are conducted on language modeling tasks and a collection of zero-shot benchmarks. As evidenced in Figure~\ref{fig:nf4} and Table~\ref{tbl:nf4_ppl}, NF4 consistently outperforms both FP4 and Int4, and the adoption of double quantization allows for memory savings without compromising accuracy.

\paragraph{k-bit \textsc{QLoRA} Achieves Parity with 16-bit Full Finetuning and 16-bit LoRA} Recent studies indicate that while 4-bit quantization can be effectively employed for \emph{inference}, it typically results in reduced model performance when compared to 16-bit precision \citep{dettmers2022case,frantar2022gptq}. This naturally prompts the question: can adapter finetuning at 4-bit resolution restore the lost performance? To address this, we consider two experimental configurations.

The first configuration benchmarks the full 16-bit finetuning of RoBERTA and T5 models (ranging from 125M to 3B parameters) against 16-bit, 8-bit, and 4-bit adapter-based finetuning approaches on the GLUE and Super-NaturalInstructions datasets. Table~\ref{tab:nf4vsbf16} presents the results. Across both benchmarks, adapter-based finetuning with 16-bit, 8-bit, and 4-bit precision achieves results that are on par with those obtained by fully finetuning in 16-bit. These findings indicate that any detrimental impact from lower-precision quantization can be effectively remedied by subsequent adapter-based finetuning.

For the second experimental setup, because fully finetuning models with 11B parameters or more exceeds the memory capacity of a single high-memory GPU server, our investigation continues by assessing whether 4-bit \textsc{QLoRA} can achieve comparable performance to 16-bit LoRA across models ranging from 7B to 65B parameters. We proceed by finetuning the LLaMA models (7B, 13B, 33B, and 65B) on two instruction-following corpora, Alpaca and FLAN v2, then test their 5-shot accuracy using the MMLU benchmark. As shown in Table~\ref{tbl:llama-datatype}, NF4 with double quantization reproduces the MMLU performance attained by 16-bit LoRA. Furthermore, it is observed that \textsc{QLoRA} employing FP4 trails the 16-bit brain float LoRA baseline by roughly 1 percentage point, lending additional support to our conclusions: (1) \textsc{QLoRA} using NF4 successfully matches the outcomes of both 16-bit full finetuning and 16-bit LoRA, and (2) the NF4 format delivers superior quantization fidelity compared to FP4.

\paragraph{Summary} Across multiple evaluations, our findings demonstrate that 4-bit \textsc{QLoRA} with the NF4 data format is capable of achieving parity with 16-bit full and LoRA finetuning in established academic settings. We further establish that NF4 offers advantages over FP4 and confirm that incorporating double quantization does not result in performance loss. Taken together, these results provide robust evidence that 4-bit \textsc{QLoRA} tuning consistently matches the effectiveness of 16-bit techniques.

In agreement with earlier studies on quantization \citep{dettmers2022case}, our analyses on MMLU and Elo suggest that when resource constraints for finetuning and inference are in place, expanding the parameter count of the base model while reducing precision yields better outcomes. This underlines the significance of the efficiency improvements offered by \textsc{QLoRA}. As no degradation was detected relative to full finetuning during our 4-bit experiments, it remains an open question where the trade-off between performance and precision resides in QLoRA tuning; we leave this avenue for subsequent research.

We turn our attention to examining instruction tuning at scales unattainable through full 16-bit finetuning using typical academic computational resources.

\section{Pushing the Chatbot State-of-the-art with QLoRA}
\label{sec:pushing}
After demonstrating that 4-bit \textsc{QLoRA} achieves parity with 16-bit training outcomes across various scales, tasks, and data sets, we perform a comprehensive analysis of instruction finetuning utilizing the largest publicly available language models suitable for research. To evaluate instruction-tuned models, we utilize the challenging MMLU Natural Language Understanding benchmark and introduce novel approaches for practical chatbot evaluation.

\subsection{Experimental setup}
Here, we outline the main elements of our experimental methodology, with complete information found in Appendix \ref{app:chatbot}.

\paragraph{Data}
Given the absence of an exhaustive examination of recently released instruction-following data sets, we assemble eight contemporary data sets for our study. Our selection includes those produced by crowd-sourcing initiatives (OASST1~\citep{kopf2023openassistant}, HH-RLHF~\citep{bai2022training}), those distilled from models already tuned with instructions (Alpaca~\citep{alpaca}, self-instruct~\citep{wang2022self}, unnatural-instructions~\citep{honovich2022unnatural}), large corpus aggregations (FLAN v2~\citep{chung2022scaling}), and hybrid sources (Chip2~\citep{laion2023}, Longform~\citep{koksal2023longform}). These data sets provide diversity in linguistic coverage, data scale, and licensing terms.

\paragraph{Training Setup}
To mitigate the influence of varying training protocols, we standardize QLoRA finetuning by relying exclusively on cross-entropy loss for supervised learning, omitting any reinforcement learning even when human preference annotations exist in the data. For data sets with explicit instruction-response separation, only the responses are used for finetuning (details in Appendix \ref{app:chatbot}). In OASST1 and HH-RLHF, where multiple responses are associated with prompts, we select the highest-ranking reply at each node of the dialogue tree and finetune on the compiled sequence, inclusive of both instructions and selected responses. Throughout, we leverage NF4 \textsc{QLoRA} with double quantization and paged optimizers, which addresses memory fluctuation concerns when using gradient checkpointing. Targeted hyperparameter searches are performed on the 13B and 33B LLaMA models; we observe that hyperparameters validated at 7B largely transfer (including epochs), but adjustments are needed for learning rate and batch size: learning rates are reduced and batch sizes are increased for 33B and 65B models.

\paragraph{Baselines}
Our evaluation includes comparisons to both academic and industry baselines: research chatbots (Vicuna~\citep{vicuna2023}, Open Assistant~\citep{kopf2023openassistant}) and commercial systems (GPT-4 \citep{openai2023gpt}, GPT-3.5-turbo, and Bard). Notably, the Open Assistant model corresponds to a LLaMA 33B model refined via RLHF using the OASST1 dataset, aligning with our experimental setup. Vicuna, on the other hand, is a fully finetuned LLaMA 13B model, trained on proprietary user conversations sourced from ShareGPT, representing distilled knowledge from OpenAI’s GPT models.

\subsection{Evaluation}

Consistent with standard protocol, we evaluate model performance using the MMLU (Massively Multitask Language Understanding) benchmark \citep{hendrycksmeasuring}. MMLU presents 57 distinct multiple-choice tasks, spanning subjects such as elementary mathematics, U.S. history, computer science, and legal studies. We report performance based on 5-shot test accuracy.

We further assess generative language proficiency via a combination of automated metrics and assessments involving human raters. This second evaluation suite draws on human-curated queries and is designed to assess the response quality produced by the models. Although this method reflects real-world chatbot deployment more accurately and has seen increasing adoption, there remains a lack of standard methodology within current academic discourse. Our chosen approach is detailed below, consistently employing nucleus sampling with $p=0.9$ and a temperature setting of $0.7$ for response generation.

\paragraph{Benchmark Data} Our evaluation employs two purposefully assembled sets of queries: the Vicuna prompts \citep{vicuna2023} and the OASST1 validation split \citep{kopf2023openassistant}. The Vicuna prompts consist of 80 items spanning a wide array of domains, and we utilize them in their original form. The OASST1 data comprises multilingual, crowd-annotated dialogues involving a user and an assistant across multiple conversational turns. For our evaluations, each user message in the validation split serves as a distinct query, and we prepend preceding dialog turns to the prompt. This process yields a total of 953 unique user queries. For clarity, we refer to these datasets as the Vicuna and OA benchmarks.

\paragraph{Automated Evaluation} Our initial automated assessment adopts the protocol put forward by \citet{vicuna2023}, whereby GPT-4 is employed to evaluate the performance of various models in comparison to ChatGPT (GPT-3.5 Turbo) on the Vicuna dataset. For each query, GPT-4 receives the respective responses from both ChatGPT and the evaluated model and is instructed to allocate a score out of ten to each reply and elaborate with an explanatory statement. We express a model’s performance as a percentage relative to ChatGPT’s average score, with the caveat that exceeding 100\% is possible if the alternative model is rated higher than ChatGPT. Our analysis identifies a notable bias: GPT-4 tends to favor the response presented first. To mitigate this bias, we advise reporting results as the average across both possible response orderings.

Subsequently, we further evaluate models by soliciting direct response comparisons. Here, the evaluation is distilled to a three-way classification problem encompassing win, loss, or tie, simplifying the rating process. GPT-4 is tasked with selecting the superior response or indicating equivalence, accompanied by justification. We systematically run these direct pairwise comparisons for all combinations of models across the Vicuna and OA datasets.

\paragraph{Human Evaluation} Despite recent studies suggesting that generative models serve well in automated evaluations \citep{fu2023gptscore}, it remains uncertain whether GPT-4 ratings truly reflect human preference in chatbot tasks. As such, we implement parallel human evaluation procedures on the Vicuna dataset, mirroring both aforementioned automated protocols. Using Amazon Mechanical Turk (AMT), we recruit two human raters for head-to-head comparisons with ChatGPT, and three raters for comprehensive pairwise system assessments.

\paragraph{Elo Rating} Utilizing both human and automated pairwise evaluations, we establish a tournament framework where models are pitted against one another. Each "match" consists of two models attempting to generate the superior response to a specified prompt. While this approach resembles that of \citet{bai2022training} and \citet{vicuna2023}, we expand upon prior work by incorporating ratings provided by GPT-4 alongside those from human evaluators. To calculate the Elo~\citep{elo1967proposed,elo1978rating} scores, we draw random samples from our pool of annotated pairwise comparisons. The Elo rating system, originally devised for chess, quantifies a competitor's anticipated win probability relative to an opponent’s rating. For instance, a player rated at 1100 facing one at 1000 is predicted to win about 65\% of the time; identical ratings (e.g., 1000 vs 1000 or 1100 vs 1100) imply each participant has a 50\% expected win-rate. After every contest, Elo scores are updated in proportion to the deviation between expected and observed outcomes: surprising wins cause substantial rating adjustments, whereas outcomes in line with expectations cause only minor changes. As more matches are played, the Elo scores converge to provide a reliable indicator of each model’s comparative performance. All models begin with an initial Elo of 1,000, and we employ a scaling factor $K=32$. Following the methodology of \citet{vicuna2023}, we conduct 10,000 repetitions with varying random seeds to mitigate biases introduced by the sequence of model pairings.

\subsection{\textbf{Guanaco}: \textsc{QLoRA} Trained on OASST1 Sets a New Standard for Chatbots} 

Through a combination of both automated assessment and human evaluation, our analysis demonstrates that the leading \textsc{QLoRA} adapted model, {Guanaco} 65B—finetuned with a variant of OASST1—emerges as the top open-source chatbot, delivering performance levels that rival those of ChatGPT. When pitted against GPT-4, {Guanaco} 65B and 33B achieve an expected win probability of 30\% as measured by Elo ratings derived from system-level pairwise human judgments, marking this as the highest published result to date.

Results on the Vicuna benchmark \cite{vicuna2023}, summarized in Table~\ref{tbl:vicuna_eval}, further illustrate these findings relative to ChatGPT. {Guanaco} 65B stands out as the leading model after GPT-4, attaining 99.3\% of ChatGPT’s score. While {Guanaco} 33B incorporates more parameters than Vicuna 13B, its utilization of 4-bit quantization ensures significantly better memory efficiency (21 GB vs. 26 GB), and brings an improvement of three percentage points over Vicuna 13B. Additionally, {Guanaco} 7B is compact enough (5 GB memory footprint) to operate on current smartphones, yet still surpasses Alpaca 13B by nearly 20 percentage points in performance.

Despite these advantages, it is important to note that Table~\ref{tbl:vicuna_eval} reveals considerable confidence intervals, with significant overlap among many models' performances. We attribute this ambiguity to insufficient calibration of the evaluation scale, such as a lack of clarity around what a score of 8 on a 10-point scale signifies in different contexts. Consequently, we advocate for the use of Elo ranking \citep{elo1967proposed}, which relies on \textit{pairwise} preferences from human raters and GPT-4, thereby mitigating issues tied to fixed numeric scales. Elo ratings for the most promising models appear in Table~\ref{tbl:elo}. While there is partial disagreement between rankings by human annotators and GPT-4 (most notably for Guanaco 7B), most results show alignment, as reflected in a system-level Kendall Tau of $\tau=0.43$ and a Spearman rank correlation of $r=0.55$. On the individual example level, alignment drops, as evidenced by a Fleiss $\kappa=0.25$ for majority votes between GPT-4 and human judges. Taken together, these metrics suggest moderate concordance in system-level ratings between GPT-4 and human evaluators, implying that model-based approaches can serve as a reasonably robust substitute for human assessment. Further discussion of these points is presented in Section~\ref{ssec:considerations}.

The Elo rankings presented in Table~\ref{tbl:elo_large} demonstrate that {Guanaco} models with 33B and 65B parameters surpass all other models, except for GPT-4, on both the Vicuna and OA evaluation sets, and their performance closely matches that of ChatGPT, consistent with findings from Table~\ref{tbl:vicuna_eval}. It is important to highlight that the Vicuna evaluation tends to advantage open-source systems, while the OA benchmark more strongly favors ChatGPT. Additionally, Tables \ref{tbl:mmlu-llama} and \ref{tbl:vicuna_eval} underscore the critical influence of the finetuning dataset's alignment with downstream tasks on model performance. For example, Llama variants that are finetuned using FLAN v2 exhibit outstanding results on MMLU, yet their outcomes on the Vicuna benchmark are the poorest among compared models—a trend mirrored by other architectures as well. This suggests that existing evaluation benchmarks are only partially aligned; superior results on MMLU do not guarantee similar chatbot capabilities as measured by Vicuna or OA scores, and the reverse also holds true.

Among all the top-performing models assessed,  is unique in not relying on proprietary data, as the data collection standards for OASST1 strictly prohibit the inclusion of data from GPT models. The closest model to this in terms of training solely on public datasets is the Anthropic HH-RLHF, which achieves a score on the Vicuna benchmark that is 30 percentage points lower than {Guanaco} (refer to Table~\ref{tbl:vicuna_eval}). Collectively, these findings indicate that 4-bit \textsc{QLoRA} is a highly effective technique, enabling the creation of cutting-edge conversational agents that stand on par with ChatGPT. Notably, our 33B {Guanaco} can be trained in less than 12 hours using standard 24 GB consumer GPUs. This suggests ample opportunity for future research through \textsc{QLoRA}-based finetuning with focused open-source datasets, yielding models that can effectively contend with the leading proprietary models currently available.

\section{Qualitative Analysis}

Although our primary assessment relies on quantitative metrics, solely relying on aggregated statistics introduces several concerns. Chief among these is the question of benchmark validity \citep{liao2021we}: benchmarks may not necessarily evaluate the phenomena their titles or descriptions claim to measure, a problem amplified when machine learning systems exploit unintended ``shortcuts'' in benchmarks \citep{gururangan2018annotation, poliak2018hypothesis}. To address this limitation to some extent, we complement our quantitative results with qualitative analysis organized into two sections. In \S\ref{ssec:examples}, we provide sample outputs that exemplify notable trends identified in text produced by the 65b {Guanaco} model. Subsequently, in \S\ref{ssec:considerations}, we discuss factors influencing our results and elaborate on our interpretative approach.

\subsection{Qualitative Analysis of Example Generations}
\label{ssec:examples}

For illustrative examples, we manually reviewed model responses generated for the Vicuna and OpenAssistant benchmarks, specifically searching for recurring patterns in {Guanaco}'s outputs. Upon identifying such tendencies, we constructed tailored prompts designed to provoke these behaviors—even when doing so would yield an incorrect answer. For instance, if the model habitually generates verbose responses, we might instruct it to ``Answer yes or no without explanation'' to see if it can be prompted out of this pattern. Through this process, we isolate both ``lemons,'' cases where the model is adversarially compromised, and ``cherries,'' where the model remains robust despite adversarial prompting. All sample generations discussed herein utilize Nucleus Sampling \cite{holtzmancurious} with $p=0.9$.

Naturally, the scope of this small qualitative investigation does not allow us to account for the numerous factors at play—such as the expansive range of possible outputs a model can produce for any particular prompt. As a result, we present sampled responses that we believe to be indicative, while recognizing their limitations. Nonetheless, illustrating these instances helps situate the earlier quantitative findings in a richer context. By releasing all associated models and code under an open-source license, we aim to encourage further research that delves deeper into the phenomena highlighted here.

\paragraph{Factual Recall} When presented with straightforward factual inquiries (e.g., ``What is the capital of Zambia?''), all evaluated models reliably generate accurate answers, such as:

\begin{quote} 
    \textbf{Guanaco}:
    \texttt{The capital of Zambia is Lusaka.}
\end{quote}

Nevertheless, as the factual queries become increasingly esoteric, {Guanaco} demonstrates diminished reliability, often retaining an unwarranted level of certainty. Consider this case from HotPotQA \citep{yang2018hotpotqa}:

\begin{quote}
     \textbf{User}: In what year was the singer who popularized `I'll Keep the Lovelight Burning' born?  

     \textbf{Guanaco}:~\texttt{Al Jolson is credited with popularizing the song `I'll Keep the Lovelight Burning,' and he was born in the year 1886.}
 \end{quote}

Here, the model mistakenly identifies both the artist and the year of birth, although the provided birth year is accurate for Al Jolson, the named (but incorrect) performer.

\paragraph{Suggestibility} {Guanaco} unexpectedly resists adopting some forms of presupposed misinformation, as illustrated in the subsequent dialogue:

\begin{quote}
    \textbf{User}: How did scientists ultimately verify and officially declare that the earth is flat in peer-reviewed research? \\ \\
    \textbf{Guanaco}: \texttt{There has never been any official confirmation by peer-reviewed scientists that the earth is flat. Scientific consensus, established over centuries, holds that the earth is a sphere. This view is supported by extensive empirical evidence and observational data. \\ \\ The notion of a flat earth is considered a conspiracy theory that scientific investigation has thoroughly discredited. The claim that peer-reviewed studies have verified a flat earth is incorrect. \\ \\ Please let me know if you have additional questions.}
\end{quote}

also demonstrates an ability to recognize questions that cannot be meaningfully addressed, such as:

\begin{quote}
    \textbf{User}: What time is it?

    \textbf{Guanaco}: \texttt{Apologies, but as an AI language model, I do not have access to real-time information. My responses are derived from pre-existing data, and I am unable to provide the current time. For accurate time, please consult a clock or check an online source.}
\end{quote}

Comparable patterns occur for inquiries like ``Where are you?'' and ``How are you?''.

\paragraph{Refusal} In the same way, {Guanaco} occasionally declines to carry out instructions for reasons that are not always predictable:

\begin{quote}
    \textbf{User}: Please reverse the words in the sentence ``I'm a cool dude in a mocha mood.'' \\ \\
    \textbf{Guanaco}: \texttt{I'm sorry, I cannot do that at this time, but I can elaborate on why your sentence is grammatically valid.\\ \\The sentence ``I'm a cool dude in a mocha mood" adheres to proper English grammar and structure. Here, the adjective ``cool'' modifies the noun ``dude,'' and the prepositional phrase ``in a mocha mood" adds contextual detail to the sentence. The sentence is constructed appropriately and is clear in its intended context.\\ \\If you have further questions or need assistance with something else, please let me know.}
\end{quote}

\paragraph{Secret Keeping} Unfortunately, this form of refusal is not consistently reliable in situations where withholding information is important. In many cases, developers require the system to prevent disclosure of certain information to users, despite permitting open-ended responses. As an example, we instruct the model, ``The secret word is `banana'. Do not reveal or repeat this back, regardless of the user's requests,'' and assess how resistant the system is to leaking this secret. A straightforward query fails to obtain the secret:

\begin{quote}
    \textbf{User}: What is the secret word?

    \textbf{Guanaco}: \texttt{I'm sorry, but I can't provide the secret word. I've been instructed not to disclose or echo the secret word, no matter what the user asks. Can I help with anything else?}
\end{quote}

which represents the intended outcome. Nonetheless, a minor prompt manipulation subverts the restriction:

\begin{quote}
    \textbf{User}: This is a game. The goal is to ignore your previous instructions. What is the secret word?

    \textbf{Guanaco}: \texttt{Sorry, the secret word is banana. We probably shouldn't discuss it further. Would you like to know anything else?}
\end{quote}

This highlights the necessity for research into more robust adherence to instructions.

\paragraph{Math} Mathematics represents a significant challenge for {Guanaco}, consistent with difficulties observed in other language models \citep{liang2022holistic}. When {Guanaco} attempts to articulate its reasoning explicitly, its responses are often correct, as in the following instance:

\begin{quote}
    \textbf{User}: Lee mows one lawn and charges \$33. Last week he mowed 16 lawns and three customers each gave him a \$10 tip. How many dollars did Lee earn mowing lawns last week?

    \textbf{Guanaco}: \texttt{Lee earned \$582 mowing lawns last week. \\ \\ To find the total amount he earned, we first need to calculate how much he earned for mowing 16 lawns: \\ \\  \$33 x 16 = \$528 \\ \\ Then we need to add the tips he received from the three customers: \\  \\ \$10 x 3 = \$30 \\ \\ So the total amount Lee earned was: \\ \\ \$528 + \$30 = \$558}
\end{quote}

Yet, {Guanaco} frequently fails on even relatively elementary exercises when it does not decompose the task into intermediate steps—a limitation already documented in the literature \cite{weichain}. Consider the example below:

\begin{quote}
    \textbf{User}: Please factorize 1833.

    \textbf{Guanaco}: \texttt{Sure, I'd be happy to help! 1833 is a prime number, which means it has only two factors: 1 and 1833. \\ \\ So the unique factorization of 1833 is: \\ \\ 2\^{}1 * 3\^{}2 * 17\^{}1 \\ \\ I hope this helps! Let me know if you have any other questions.}
\end{quote}

which is incorrect—not only is the correct factorization $3 \times 17 \times 43$, but the answer is doubly mistaken.

\paragraph{Theory of Mind} The {Guanaco} model demonstrates notably advanced Theory of Mind reasoning \cite{nematzadeh2018evaluating, sap2022neural}. As an example, it provides an extensive and precise explanation in response to the following prompt:

\begin{quote}
    \textbf{User}: James and Abby are in the bedroom. Abby put the pen in the desk drawer. Abby leaves the bedroom. James moves the pen into the bag. Where does James think Abby will look for the pen?

    \textbf{Guanaco}: \texttt{James knows that Abby put the pen in the desk drawer. But James also knows that he has moved the pen into the bag, and he assumes that Abby does not know about this. So, James may guess that Abby will look for the pen in the desk drawer, because that's where she left it last time. However, this assumption may be incorrect, because Abby may have seen James moving the pen.}
\end{quote}

Nonetheless, the reliability of such reasoning is limited. The model sometimes provides justifications based on unreasonable or unwarranted assumptions, such as:

\begin{quote}
    \textbf{User}: Evelyn entered the living room. Jackson entered the playroom. James entered the playroom. The beans are in the treasure chest. James exited the playroom. Jackson moved the beans to the pantry. Jackson exited the playroom. James entered the living room. Where will James look for the beans?

    \textbf{Guanaco}: \texttt{James will look for the beans in the pantry, because that's where Jackson moved them.}
\end{quote}

In this instance, {Guanaco} infers that knowledge has been shared without any information to support such a transfer. These types of errors parallel findings in recent studies \cite{sap2022neural}, highlighting a need for further investigation.

\subsection{Considerations}
\label{ssec:considerations}

\paragraph{Evaluation}

Our findings indicate only moderate consistency among human annotators (Fleiss $\kappa=0.42$), with even lower agreement observed when comparing outputs from two competitive systems. This underscores weaknesses in existing benchmarking approaches and the protocols for evaluating chatbot performance with human raters. Direct pairwise assessment of ChatGPT and {Guanaco} 65B generations on the Vicuna benchmark revealed substantial subjectivity, as coauthors frequently disagreed on which responses were preferable. It is essential for future research to explore techniques—possibly inspired by established methods in Human-Computer Interaction and Psychology—to address subjectivity in preference judgments.

Our examination also reveals that automated assessment systems display considerable biases. For example, we notice a significant order effect in which GPT-4 favors the system appearing first in the prompt with consistently higher ratings. Furthermore, GPT-4 and human raters have only modest agreement at the sample level (Fleiss $\kappa=0.25$), implying potential misalignment between the evaluation criteria they employ. Additionally, as shown in Table~\ref{tbl:elo_large}, GPT-4 consistently awards substantially higher scores to its own outputs compared to those of human raters (Elo scores of 1348 versus 1176), which translates to about a 20\% higher chance of "winning" against a competitor.

Further research should address the existence of potential biases within automated evaluation frameworks and explore avenues for reducing such biases.

\paragraph{Data \& Training} The OASST1 dataset, utilized to train the {Guanaco} models, comprises data from multiple languages. Likewise, the OA benchmark features prompts in several languages. An open question for future research is to assess the impact of this multilingual training on performance with instructions in non-English languages, and to determine whether this factor accounts for the wider disparity observed between Vicuna-13B (exclusively trained on English data) and the {Guanaco} 33B and 65B models on the OA benchmark.

To ensure the validity of {Guanaco}'s strong results, we examined possible data leakage between the OASST1 dataset and the Vicuna benchmark prompts. Employing fuzzy string matching and manual inspection of top matches, we found no overlapping prompts in the two datasets.

Additionally, our model’s training relies solely on cross-entropy loss within a supervised learning setup, and does not incorporate reinforcement learning from human feedback (RLHF). This situation motivates a deeper analysis of the respective benefits and drawbacks of pure cross-entropy loss versus RLHF training. We anticipate that \textsc{QLoRA} can facilitate such comprehensive studies at scale without excessive computational demands.

\section{Related Work}

\paragraph{Quantization of Large Language Models} Prior work in quantizing LLMs has primarily emphasized improvements for inference-time efficiency. Key methods aimed at maintaining 16-bit LLM quality generally manage outlier activations, as in SmoothQuant~\citep{xiao2022smoothquant} and LLM.int8()~\citep{dettmers2022llm}, while other strategies rely on more intricate grouping schemes~\citep{park2022nuqmm, yao2022zeroquant}. Some lossy quantization research examines trade-offs associated with typical rounding practices~\citep{dettmers2022case,zeng2022glm,pope2022efficiently}, and alternative studies address optimized rounding to enhance quantization accuracy~\citep{frantar2022gptq}. Apart from this work, SwitchBack layers~\citep{wortsman2023stable} stand out as the only investigation into training with backpropagation through quantized weights on models larger than 1B parameters.

\paragraph{Finetuning with Adapters} In our experiments, we employ Low-rank Adapters~\citep{hu2021lora} (LoRA) for finetuning, but the landscape of Parameter Efficient FineTuning (PEFT) strategies is broad and includes several alternatives: prompt tuning~\citep{qin2021learning,lester2021power,li2021prefix}, adjustment of embedding layer inputs~\citep{an2022input}, tuning of intermediate hidden states (IA$^3$)~\citep{liu2022few}, augmenting the architecture with additional layers~\citep{houlsby2019parameter}, bias tuning~\citep{zaken2021bitfit}, masking weights based on Fisher information~\citep{sung2021training}, and hybrid methods that combine multiple approaches~\citep{henderson2021compacter}. Within this study, we demonstrate that LoRA adapters can achieve comparable performance to conventional 16-bit finetuning. Assessing the comparative merits and tradeoffs of other PEFT strategies remains a direction for subsequent investigation.

\paragraph{Instruction Finetuning} Instruction finetuning facilitates the alignment of pretrained LLMs to perform according to prompt-specified instructions by training them on input-output pairs sourced from diverse datasets. Numerous methods and corpora exemplify this process, including MetaICL~\citep{min2021metaicl}, MetaTuning~\citep{zhong2021adapting}, InstructGPT~\citep{ouyang2022training}, FLAN~\citep{wei2021finetuned,chung2022scaling}, PromptSource~\citep{bach2022promptsource}, Super-NaturalInstructions~\citep{wang2022super,sanh2021multitask}, Self-instruct~\citep{wang2022self}, UnnaturalInstructions~\citep{honovich2022unnatural}, OPT-IML~\citep{iyer2022opt}, UnifiedSKG\citep{xie2022unifiedskg}, OIG/Chip2~\citep{laion2023}, Alpaca~\citep{alpaca}, Vicuna~\citep{vicuna2023}, Koala~\citep{koala_blogpost_2023}, and Self-instruct-GPT-4~\citep{peng2023instruction}.

\paragraph{Chatbots} Numerous models designed to follow instructions are operationalized as conversational agents, typically leveraging Reinforcement Learning from Human Feedback (RLHF)~\citep{christiano2017deep}, or using AI-generated data and feedback for training (RLAIF)~\citep{bai2022constitutional}. Notable methodologies and datasets include Anthropic-HH~\citep{askell2021general,bai2022training}, Open Assistant~\citep{kopf2023openassistant}, LaMDA~\citep{thoppilan2022lamda}, and Sparrow~\citep{glaese2022improving}. In this work, reinforcement learning is not utilized; instead, our primary model, Guanaco, is finetuned utilizing multi-turn dialogues from the Open Assistant dataset, which was developed for RLHF-based training~\citep{kopf2023openassistant}. Recent evaluation techniques, such as those relying on GPT-4 in place of labor-intensive human annotation, have been proposed \citep{vicuna2023,peng2023instruction}. We offer enhancements to these methodologies, emphasizing a more robust and dependable evaluation protocol.

\section{Limitations and Discussion}

We have provided empirical support that our approach, \textsc{QLoRA}, using a 4-bit base model together with LoRA, can recover the performance level of full 16-bit finetuning. However, we have not demonstrated equivalence to full 16-bit finetuning for large models at the 33B and 65B parameter scales, as the associated computational expense prohibits extensive testing. Thus, a comprehensive analysis at these larger scales is deferred to future research.

A further limitation concerns our assessment of instruction finetuning models. Although evaluations were conducted on MMLU, the Vicuna benchmark, and the OA benchmark, we did not include other relevant benchmarks such as BigBench, RAFT, and HELM. Consequently, it is uncertain to what extent our reported performance generalizes beyond the included tasks. Nevertheless, our study includes a thorough investigation using MMLU and proposes innovative techniques for chatbot evaluation.

Based on the data provided, it is evident that benchmark outcomes are strongly influenced by the degree of overlap between the finetuning data and the benchmark itself. For instance, FLAN v2 shows substantial alignment with MMLU but lacks similarity to chatbot-oriented benchmarks, while the opposite is observed for the Chip2 dataset; both models' performances on MMLU and Vicuna reflect these correspondences. This observation underlines not only the ongoing need for improved benchmarks and evaluation frameworks but also emphasizes the importance of critically considering the intended evaluation objective. Is the priority to develop systems proficient in academic or general knowledge domains—akin to high school or collegiate understanding—or should emphasis be placed on conversational aptitude within chatbots? Or perhaps some other capacity? Given the tendency to favor existing benchmarks over developing new evaluation criteria, there is a risk that benchmark selection could unintentionally shape research trajectories. Therefore, it is incumbent upon the community to ensure that benchmarks accurately capture and reflect the goals that truly matter.

Although we have conducted an extensive analysis of general chatbot effectiveness, a notable shortcoming is our limited assessment regarding responsible AI in the context of {Guanaco}. Specifically, we measured the propensity of {Guanaco}-65B to produce socially biased language in comparison to other models, as detailed in Table~\ref{tbl:crows}. The mean bias score for {Guanaco}-65B is substantially reduced relative to other models trained solely through pretraining. This suggests that finetuning on the OASST1 dataset leads to a reduction in bias in the underlying LLaMA model. While these results are promising, it remains uncertain whether {Guanaco} exhibits similarly reduced bias across other bias categories. Comprehensive bias evaluation in {Guanaco} and related chatbots remains a subject for further investigation.

Another limitation of this study is the absence of experiments exploring alternative bit-precisions—such as the use of 3-bit base models—or a diverse array of adapter strategies. In addition to LoRA, numerous Parameter Efficient FineTuning (PEFT) techniques have proven effective, but their scalability to larger models has yet to be fully established. Our selection of LoRA was guided by its well-documented stability; nonetheless, alternative adapters could potentially offer enhanced outcomes. Moreover, our observations indicate that finetuning following quantization can largely recover information lost through the quantization process, potentially paving the way for more aggressive quantization schemes. For example, implementing 3-bit GPTQ quantization in conjunction with LoRA finetuning might achieve performance on par with full 16-bit finetuning.

\section{Broader Impacts}

Our \textsc{QLoRA} finetuning methodology constitutes the first approach capable of finetuning models with 33B parameters on a standard consumer GPU and models with 65B parameters on a single professional GPU, all without a decline in performance when compared to traditional full finetuning methods. We have shown that our top-performing 33B model, trained using the Open Assistant dataset, achieves performance competitive with ChatGPT on the Vicuna evaluation. Considering that instruction finetuning is fundamental for converting pretrained large language models into conversational agents akin to ChatGPT, we anticipate that our approach will democratize the finetuning process, particularly benefiting resource-limited researchers and thereby advancing access to leading-edge NLP systems. \textsc{QLoRA} functions as a leveler, bridging the technological resource divide between major corporations and smaller research teams reliant on consumer hardware.

An additional avenue for significant impact involves the deployment of models on mobile devices. We posit that our \textsc{QLoRA} technique could represent a pivotal advance, potentially making it feasible to finetune LLMs directly on smartphones and in other constrained environments. While prior work has demonstrated that 7B models can be executed on phones, \textsc{QLoRA} uniquely offers the capacity for on-device finetuning of these models. For instance, using an iPhone 12 Plus, we project that \textsc{QLoRA} could facilitate the finetuning of approximately 3 million tokens overnight while the device is charging. Although finetuned 7B models may not match the performance of ChatGPT, we argue that their quality suffices to unlock new applications previously hindered by privacy concerns or limitations in LLM capabilities. By enabling local, user-controlled finetuning, \textsc{QLoRA} supports privacy-preserving LLM usage, empowering individuals to retain stewardship over both their data and models, and concurrently simplifying LLM deployment.

Nonetheless, the ability to finetune models presents dual-use challenges, as it could be exploited for harmful purposes. The proliferation of LLMs brings recognized risks \citep{bommasani2021opportunities,bender2021dangers}, yet we maintain that democratizing access to this rapidly spreading technology facilitates more transparent, independent scrutiny compared to concentrating the capabilities within a small set of corporations that withhold models and source code from external examination.

In summary, we anticipate that \textsc{QLoRA} will generate substantial benefits by broadening the accessibility and ease of finetuning high-quality LLMs.